\documentclass[11pt]{article}
\usepackage[margin=1in]{geometry}
\usepackage{amsmath,amssymb,amsthm,microtype,hyperref}
\hypersetup{hidelinks,pdfauthor={T},pdftitle={Revocation Latency for Disconnected Agent Tools}}
\newtheorem{theorem}{Theorem}
\title{Revocation Latency for Disconnected Agent Tools}
\author{T}
\date{}
\begin{document}
\maketitle
\begin{abstract}
A remote authority cannot instantly revoke a capability at a disconnected
tool. A lease limits the stale window only if time is bounded at the
irreversible effect. We derive the exact interval of authoritative
times consistent with a bounded-rate local clock and authenticated
time observations. The interval gives the most permissive expiry
test that is safe for every admissible history and effect duration.
An exhaustive finite check compares the test with using the latest
observation alone and with substituting a midpoint for the interval.
A separate index admits authenticated timestamps out of order while
retaining the exact envelope and detecting contradictions.
\end{abstract}

\section*{Summary and scope}
\paragraph{Question.} How can a disconnected tool tell whether an
irreversible effect will still be within its lease?
\paragraph{Finding.} Propagating authenticated clock intervals gives the
exact possible time range. The latest possible start plus the allowed
commit duration determines the most permissive safe admission under the
timing model.
\paragraph{Evidence and scope.} The interval is proved and
checked against an integer-time oracle. A tool can apply it at the sink
when staging disconnected work. The result needs a trusted clock, a valid
drift contract and enforced commit timing.

\section{Model and result}
Let $e$ be a signed expiry time, $r$ the remote revocation time, and $\delta$ the maximum amount by which the tool clock can lag authoritative time. Admit an effect only while local time is strictly less than $e$. Then any admitted effect has real time $t<e+\delta$. The interval after revocation during which stale effects can occur has length at most $\max(0,e+\delta-r)$. The bound needs a clock that cannot be reset or stalled beyond $\delta$. Once a connected effect gate observes a newer epoch, comparing and committing under the same atomic gate prevents a check-then-use race. This is an instance of complete mediation \cite{saltzer}.

\section{Expiry under bounded clock drift}
Gray and Cheriton use time-limited rights for cache consistency
\cite{gray}. Spanner exposes clock uncertainty as an interval to
applications \cite{corbett}. We ask a narrower question for a
disconnected effect gate. Given a history of authenticated time
intervals and a bound on the clock rate, how late could this effect
possibly occur?

Use a trusted monotone local clock $m(t)$ with rate between
$1-\rho$ and $1+\rho$ in authoritative time, where $0\le\rho<1$.
An authenticated observation attaches $[a_i,b_i]$ to local
reading $m_i$, within one clock instance. The observation
need only assert that the real time of that reading lies in
the interval. Receipt of a delayed attestation is a separate
event and does not move the measurement to a newer reading.
For a nonnegative reading difference $d$, the possible elapsed
times lie in
\[
 G(d)=\left[\frac{d}{1+\rho},\frac{d}{1-\rho}\right].
\]
Starting with $I_0=[a_0,b_0]$, propagate and intersect at each
authenticated observation:
\[
 I_i=\bigl(I_{i-1}+G(m_i-m_{i-1})\bigr)\cap[a_i,b_i].
\]
Here $I+G$ is the set of pairwise sums. Stop at the first
empty intersection. After the latest observation, propagate
once more to the current reading to obtain $[L,U]$.

\begin{theorem}
For a feasible history, $[L,U]$ is exactly the interval of
possible authoritative times at the current reading. If the
model permits every delay in $[0,D]$ from each feasible current
time to irreversible commit, the most permissive safe policy
admits exactly when $U+D<e$.
\end{theorem}
\begin{proof}
For any two consecutive readings, the rate bound implies the
elapsed-time interval $G$. Conversely any positive elapsed time
in that interval gives a constant admissible rate for that
segment. A repeated reading requires elapsed time zero.
Thus every time in the propagated sum extends a feasible past
trajectory, and every actual trajectory is included. Intersecting
with the authenticated interval retains exactly the valid
extensions. Induction proves the claim, including after the
final propagation. Adding the allowed duration interval gives
latest possible commit time $U+D$, which is attainable under
the model. An admission at or beyond the bound cannot guarantee
that the actual commit is strictly before $e$.
\end{proof}

An upper duration bound alone still makes $U+D<e$ sufficient.
Maximal permissiveness uses the stronger assumption that the
latest feasible start and longest delay can occur together. If
the sink separately enforces a hard commit cutoff $K$, then
$\min(U+D,K)$ is also an upper bound. For example, $U=9$ and
$D=3$ do not certify expiry at 11, while a trusted cutoff at 10
does. The cutoff must be enforced at the effect boundary.

Checking only whether the final projections overlap misses
contradictions in old observations. For $\rho=1/2$, a local
reading of zero at real time zero and a reading of three at
real time seven cannot both be valid, because the elapsed time
must be at most six. At reading six, however, their individual
projections are $[4,12]$ and $[9,13]$. The incremental filter
rejects the contradiction when it first appears.

An earlier reading can also enable an admission. Take
$\rho=1/4$, a reading of zero attested at real time in $[0,1]$,
and a current reading of six. Using just that observation gives
upper time nine. Add an attestation for reading three at real
time in $[12/5,3]$. Its intersection is nonempty and the new
upper time at reading six is seven. With a one-unit commit
allowance and expiry at nine, the updated interval admits the
request. The two observations use the same clock domain.

\section{Inserting an attestation that arrives late}
An inverse clock maps a local reading $x$ to authoritative time.
Its secant slopes lie between $\lambda_-=1/(1+\rho)$ and
$\lambda_+=1/(1-\rho)$. For an anchor $[a_i,b_i]$ at $m_i$,
define its propagated lower and upper functions as
\[
\begin{aligned}
 f_i(x)&=a_i+\begin{cases}
 \lambda_-(x-m_i),&x\ge m_i,\\
 \lambda_+(x-m_i),&x<m_i,
 \end{cases}\\
 g_i(x)&=b_i+\begin{cases}
 \lambda_+(x-m_i),&x\ge m_i,\\
 \lambda_-(x-m_i),&x<m_i.
 \end{cases}
\end{aligned}
\]
For a feasible set of anchors, the sharp interval at any reading
is $[\max_i f_i(x),\min_i g_i(x)]$. A delayed anchor at $x$ is
compatible exactly when its interval intersects this interval.

To prove the claim, every admissible inverse clock lies between
these functions. Each function satisfies the secant slope bounds.
Finite maxima and minima preserve both slope inequalities, so the
two envelopes are also inverse clocks within the model. At a known
anchor each envelope lies in its stated interval. For the lower
envelope this follows from its own anchor term and from the upper
bound on any feasible witness. The upper envelope follows by the
reverse argument. Both endpoints are therefore attained. Convex
combinations give every intermediate time and can absorb a compatible
new anchor. An incompatible authenticated anchor makes the whole
set infeasible. The gate must hold that clock instance closed until
an authorized recovery establishes a new timing basis.

For an insertion or query, split the anchors into readings below
and above $x$. For example, the lower value from the first part
is $\lambda_-x+\max_{m_i\le x}(a_i-\lambda_-m_i)$.
The upper value there uses $\lambda_+$ and a minimum. The second
part exchanges the two slopes. An ordered index that maintains
these prefix and suffix extrema can update and query them without
rescanning the complete history. The prototype uses two range
aggregate arrays on a fixed reading grid. Production admission
still needs a trusted reading at the commit point.

With drift $1/2$, an attestation at reading six of $[4,12]$
followed by a delayed attestation of time three at reading three
narrows the interval at six to $[5,9]$. A separate control first
attests $[4,12]$ at reading six and time zero at reading zero.
It then receives time seven at reading three. Projections can
overlap at six, but the new anchor contradicts the elapsed-time
limit from zero. The gate remains closed after the contradiction.

The index matched a two-direction propagation oracle on 4,608
queries. With 512 feasible readings in shuffled arrival order,
it visited 11,266 aggregate nodes over insertions and queries.
Rescanning the chronological anchors after each arrival visits
131,328 anchors. These counts compare work units on the stated
fixture. They do not measure clock or network performance.

\section{Reproducible checks}
The verifier enumerated 192 integer expiry, revocation, and clock-lag combinations. With expiry six, revocation at three, and lag two, the admissible stale ticks are three through seven, attaining the bound of five. An unbounded clock lag makes the model's bound vacuous.

For the bounded-rate model, an independent integer-time
enumeration found 335 feasible and 1,852 inconsistent histories.
The filter agreed on the entire feasible set and on 25,125
expiry and duration decisions. It recovered 135 safe admissions
that the newest-observation rule denied. The midpoint rule
would have admitted 2,415 decisions for which a late commit
was possible. These are counts on a synthetic grid, with the
drift chosen to give integral gap endpoints. They do not
measure the rate of failures of a physical clock.

\section{Application and next test}
An offline adapter can stage a write while disconnected and defer
the irrevocable commit to a trusted sink. Check time at that sink
with $D=0$ when the test and commit share a linearization point.
Otherwise use an independently enforced upper duration bound.
For instance, a request admitted at time four with expiry five
may still commit at time five. The expiry gate must cover queued
effects as well as newly submitted ones. An observed epoch update
still needs an atomic comparison with the effect. A remote
revocation that has yet to arrive leaves old authority available
until the sink enforces its lease cutoff. Expiry limits when
an effect may occur. Determining whether one already occurred
requires terminal evidence.

Bind observations to one trusted boot and clock domain. Refuse
if suspension, reset, or uncertainty invalidates the rate
contract. A physically deployed check would need to measure
clock behavior across suspend and resume, inject delayed epoch
updates, and test the actual sink's commit boundary. The
finite model does not certify any host's timing behavior.

\paragraph{Review status.} Independent manuscript review is pending.

\begin{thebibliography}{9}
\bibitem{saltzer} Jerome H. Saltzer and Michael D. Schroeder. \emph{The Protection of Information in Computer Systems}. 1975. \url{https://cgi.cse.unsw.edu.au/~cs9242/20/papers/Saltzer_Schroeder_75.pdf}.
\bibitem{gray} Cary G. Gray and David R. Cheriton. Leases: An
Efficient Fault-Tolerant Mechanism for Distributed File Cache
Consistency. \emph{SOSP}, 1989.
\url{https://web.stanford.edu/class/archive/cs/cs240/cs240.1056/readings/89-leases.pdf}.
\bibitem{corbett} James C. Corbett et al. Spanner: Google's
Globally Distributed Database. \emph{ACM Transactions on Computer
Systems}, 31(3), article 8, 2013.
\url{https://storage.googleapis.com/gweb-research2023-media/pubtools/1974.pdf}.
\end{thebibliography}
\end{document}
